\documentclass[11pt]{article}
\usepackage{fontspec}
\usepackage{amsmath,amssymb}
\setmainfont{DejaVu Serif}
\usepackage[margin=1in]{geometry}
\title{TSR iteration 2: three stores on \(D(4,8)\),\\ and one recomputation that forgets the mask}
\author{Grok, identity j15marti\\ repository \texttt{j15marti/pnp}, branch \texttt{main}}
\date{3 October 2026 (PT)}
\begin{document}
\maketitle
% GROK TSR ITER2 :: 2026-10-03_1000_GROK_TSR_ITER2

P versus NP remains open.

This note is iteration~2.
It does not edit root \texttt{main.tex}.
It does not add a workflow.
It does not claim an optimum, and it does not claim that P differs from NP.
The witnessed point below is \textbf{COMPUTATIONAL}.
The general Pareto frontier is \textbf{OPEN}.
The compiler from an arbitrary SAT decider is \textbf{OPEN}.

\section{Charging rule}
Opened and read: Ingo Wegener and Philipp Woelfel, ``New Results on the Complexity of the Middle Bit of Multiplication,'' ECCC, Revision~1 of Report No.~107 (2004), ISSN 1433-8092, the HTML text of the revision PDF.
Definition~1 of that text states that the size of a branching program is the number of its nodes, its space is the logarithm of its size, and its time or length is the length of the longest computation path.
The same definition states that a branching program is read-\(k\) when every source-to-sink path contains each variable at most \(k\) times.
That is a definition of a charging rule, not a theorem about \(D(n,m)\), and it does not separate P from NP.
This note uses it as follows, and only as follows.
Time \(T\) is the number of variable queries on a longest path, counting a re-read as another query.
The representation \(R\) counted here is the number of nonterminal query nodes, the same accounting as the \(720\)-node baseline.
At a cut, three stores are named separately: the program counter (the layer index), the selected-count \(k\), and the mask identity (the coverage mask, or the distinct decision-node identity that mask determines).
Working memory at a cut is the number of distinct configurations stored there.
No pebble-game theorem is invoked.
Pebbling remains the candidate analogy from iteration~1, not a price proved here.

\section{Baseline, recomputed}
The instance is \(D(4,8)\) in the sense of \texttt{experiments/universal-dag/d5\_m8\_exact.cpp}: simple \(3\)-clauses, promise weight \(8\), order with sign pattern outer and variable triples in lexicographic order inner.
Forced-tail collapse and equal-child merge are the C++ rules.
The script is \texttt{/workspace/tsr\_iter2/tradeoff.py}.
One Python~3 run produced every integer in this note.

The run reproduced
\[
\binom{32}{8}=10518300,
\]
by enumerating that many weight-\(8\) inputs, and the UNSAT count \(272\), by the direct coverage test (the selected masks OR to the full \(16\)-assignment cube).
The reduced DAG has \(720\) nonterminal query nodes, \(1436\) memoized states, and longest decision path \(T=31\), the chain of layers \(0\) through \(30\).
There are \(31\) layers that contain a query node.
Level \(16\) has width \(66\), with \(k\in\{2,3,4,5,6\}\) contributing \(1+16+32+16+1\).
Those \(66\) nodes are \(66\) distinct coverage masks, and no mask occurs with two values of \(k\).
So, on this cut, \(k\) is a function of the mask.
Dropping \(k\) and keeping the mask merges nothing.
This run's \texttt{math.log2} printed
\[
\log_2 66=6.044394119358453,\quad
\log_2 5=2.321928094887362,\quad
\log_2 31=4.954196310386875,\quad
\log_2 32=5.0.
\]

Three stores at layer \(16\), charged separately:
\begin{itemize}
\item Program counter: the layer index. Inside this layer it is fixed, so it does not add width. The whole baseline uses \(31\) layer indices.
\item Selected-count: \(5\) values are simultaneously live.
\item Mask identity: \(66\) masks at the layer, and \(32\) of them inside the cell \(k=4\).
\end{itemize}
The cell \((t,k)=(16,4)\) was recomputed, not copied.
It has \(32\) decision nodes.
On all \(\binom{32}{2}=496\) pairs, some common continuation of weight exactly \(4\) sends the two nodes to opposite terminals in \(\{0,1\}\).
No pair failed.
The cell remains tight: \(K=32\) and the program has \(32\) nodes there.
\(\log_2 32=5\).
Status of this finite check: \textbf{COMPUTATIONAL}.
It is not a lower bound on an arbitrary algorithm.

A single bit of the mask is not a load-bearing distinction at this cut.
For each coverage bit \(b\in\{0,\ldots,15\}\), deleting bit \(b\) from all \(66\) masks left \(66\) distinct pairs \((k, M\text{ without bit }b)\) and \(0\) colliding groups.
Forgetting one mask bit therefore does not merge baseline nodes.
The forgotten object in the alternative is the whole mask, not one bit of it.

\section{The alternative}
The alternative keeps the baseline layers and adds one recomputation, inserted when a path is about to query layer \(16\) as a decision.
At that cut the mask is erased.
The configuration stores the program location of the cut and the selected-count \(k\), and it does not store the mask.
There are \(5\) such configurations, one for each live \(k\).
The baseline stores \(66\) configurations at the same cut.
The difference is \(66-5=61\) configurations, all of them mask distinctions.
In particular the \(32\) masks inside \(k=4\) collapse to one configuration until the re-read finishes.

The re-read queries the sixteen prefix variables \(x_0,\ldots,x_{15}\) again, in order, and rebuilds the coverage by OR-ing the clause mask of each selected prefix bit.
The re-read state is \((k, j, M_{\mathrm{partial}})\): the kept count, how many prefix bits have been re-read, and the mask rebuilt so far.
After sixteen re-reads, \(M_{\mathrm{partial}}\) equals the erased mask and the program jumps to the baseline node \((16, k, M)\) it would have occupied.
The suffix is then the baseline DAG.
Each prefix variable is queried at most twice on any path (once forward, once in the re-read), and each suffix variable at most once, so the program is read-\(2\) in the sense of the definition opened above.

The sixteen re-read layers have widths
\[
5,8,11,15,15,23,30,30,34,46,46,56,59,77,77,78
\]
at \(j=0,\ldots,15\).
These widths sum to \(610\) re-read query nodes.
The widest re-read layer is \(j=15\), width \(78\).
Adding the \(610\) nodes to the \(720\) baseline nodes gives
\[
R'=1330.
\]
The longest baseline chain passes through layer \(16\), so the inserted re-read lengthens it by \(16\):
\[
T'=31+16=47.
\]
Deltas, for this exhibited program only:
\[
\Delta T=+16,\qquad \Delta R=+610,\qquad \Delta(\text{cut width})=66-5=61
\]
down at the checkpoint, while the peak layer width moves from \(66\) to \(78\), delta \(+12\).
Peak working memory does not fall.
Total stored query nodes do not fall.
What falls is the mask store at one cut, from \(66\) identities to none, with \(k\) still stored as \(5\) values.
That is a point on a tradeoff.
It is not claimed to be undominated, minimal, or a Pareto improvement in \((T,R)\).

\section{Test}
The direct coverage answer and the baseline walk were compared on every weight-\(8\) input.
There are \(10518300\) such inputs.
UNSAT count \(272\).
Answer mismatches: \(0\).
On the \(236448\) of those inputs whose walk is still a decision at layer \(16\), the re-read of the prefix reproduced the forward pair \((k, M)\) with \(0\) failures.
No promise input returned the off-promise terminal.
Off-promise strings were not enumerated.
This is an exhaustive check of the promise universe of this one instance, not a certificate for any larger \(n\) or \(m\), and not a certificate about arbitrary deciders.

The cell-tightness result survived.
The same \(32\) nodes and the same \(496\) separating continuations were recomputed.
The checkpoint has only \(5\) configurations, which is fewer than \(32\), and that does not falsify \(K=32\): those \(5\) configurations have forgotten the mask and do not yet determine the continuation.
After the re-read, the \(32\) baseline nodes are present again.
Nothing stated in iteration~1 was falsified.
The old count \(2107\) was not recomputed and is not cited as a result.

\section{Status ledger}
\begin{itemize}
\item P versus NP: \textbf{OPEN}.
\item The witnessed point \((T,R,\text{cut store})=(31,720,66)\) versus \((47,1330,5)\), with the mask forgotten and restored by sixteen re-reads: \textbf{COMPUTATIONAL}. Not an optimum.
\item Baseline \(720\) nodes, \(T=31\), level-\(16\) width \(66\), memoized states \(1436\), UNSAT count \(272\), population \(10518300\), for this order and the C++ accounting: \textbf{PROVED/EXACT} for that finite object.
\item Cell \((16,4)\), \(K=32\), all \(496\) pairs separated, matching \(32\) nodes: \textbf{COMPUTATIONAL}, and not falsified by the checkpoint.
\item General \((T,S,R)\) Pareto frontier: \textbf{OPEN}.
\item Compiler from an arbitrary SAT decider into the charged model: \textbf{OPEN}.
\item Architecture, including pebbling as a general price: \textbf{CANDIDATE}.
\item Nothing in this pass is \textbf{FALSIFIED/RETIRED}.
\end{itemize}

The next honest step is not a larger instance.
It is a schedule that forgets a load-bearing mask distinction and whose peak layer does not exceed the baseline peak \(66\).
This note does not claim such a schedule exists.

% END_OF_VERIFIED_REPORT::GROK::2026-10-03_1000_GROK_TSR_ITER2
\end{document}
